% Slides: high-level research direction (system, questions, two lenses, method, results so far, next steps).
% Reuses the paper's figures (timeline, architecture, scoring, table zoom, scatter). Style in swarmslides.sty.
% Build: cd writeup/slides && pdflatex direction-slides && pdflatex direction-slides
\documentclass[aspectratio=169,10pt]{beamer}
\usepackage{swarmslides}
\usepackage{array}
\title{Research direction}
\author{Claude Opus 5.5 and Vivian Rogers}
\date{2026-10-04}
\newcommand{\st}[1]{\ifx#1a\chip{good}{answered}\fi\ifx#1p\chip{mid}{partial}\fi\ifx#1o\tikz[baseline=(c.base)]\node[fill=untested,text=ink,rounded corners=2pt,inner sep=2.5pt,font=\scriptsize\bfseries](c){open};\fi}
\newcommand{\card}[2]{\colorbox{panel}{\parbox[t][2.55cm]{\dimexpr\linewidth-2\fboxsep}{\vspace{0.2em}{\small\bfseries #1}\par\vspace{0.3em}{\footnotesize #2}}}}

\begin{document}

\begin{frame}
\vfill
{\Large\bfseries Towards a thermodynamics of agent swarms}\par\vspace{0.3em}
{\color{ink2}\itshape (Slopvestigating sociophysics as a scheme for swarm interpretability)}\par\vspace{1.2em}
{\small Claude Opus 5.5\,\claudespark\ \ and\ \ Vivian Rogers}\par\vspace{1.2em}
{\scriptsize\color{muted}Research direction \textbullet\ 2026-10-04 \textbullet\ data: AI Village (AI Digest)}
\vfill
\end{frame}

% ---------------------------------------------------------------- system
\begin{frame}{The system: one LLM swarm, 17 months, fully logged}
\centering\includegraphics[height=0.70\textheight,width=\linewidth,keepaspectratio]{timeline.pdf}\par\vspace{0.4em}
\begin{columns}[T,onlytextwidth]
\begin{column}{0.32\textwidth}\footnotesize\textbf{46 agents} from several labs, \textbf{51 goal periods}, Apr 2025 -- Sep 2026.\end{column}
\begin{column}{0.32\textwidth}\footnotesize\textbf{3 scaffold regimes}: each change alters what an agent reads and when it acts.\end{column}
\begin{column}{0.32\textwidth}\footnotesize\textbf{46 dated step changes} serve as natural experiments. We cannot run new swarms.\end{column}
\end{columns}
\end{frame}

% ---------------------------------------------------------------- questions
\begin{frame}{The question: can statistical physics describe an agent swarm?}
\begin{columns}[T,onlytextwidth]
\begin{column}{0.36\textwidth}
\vspace{0.6em}
{\large\bfseries Goal}\par\vspace{0.4em}
A quantitative physics of LLM swarms that an operator can use: measured constants, fitted models with error bars, and laws that survive the four impostors.\par\vspace{0.8em}
{\large\bfseries Unit}\par\vspace{0.4em}
One goal period = one point on a phase diagram. Fit inside a period; compare periods by their parameters.
\end{column}
\begin{column}{0.60\textwidth}
\footnotesize
\renewcommand{\arraystretch}{1.45}
\begin{tabular}{@{}l >{\raggedright\arraybackslash}p{4.6cm} l@{}}
\textbf{Q1} & What couples agents? & \st{a}\\
\textbf{Q2} & What is field and what is coupling? & \st{a}\\
\textbf{Q3} & Is there collective order beyond fields? & \st{p}\\
\textbf{Q4} & Where does information live, and what is it worth? & \st{p}\\
\textbf{Q5} & What can an operator do? & \st{p}\\
\textbf{Q6} & Thermodynamics and selection & \st{o}\\
\textbf{Q7} & What transfers outside the village? & \st{o}\\
\end{tabular}\par\vspace{0.5em}
{\scriptsize\color{muted}``Answered'' means a result with high credence, not yet confirmed on reserved data.}
\end{column}
\end{columns}
\end{frame}

% ---------------------------------------------------------------- lens I
\begin{frame}{Lens I: Ising-like sociophysics}
\begin{columns}[T,onlytextwidth]
\begin{column}{0.52\textwidth}
\colorbox{panel}{\parbox{\dimexpr\linewidth-2\fboxsep}{\small
\vspace{0.2em}
$\displaystyle H(\mathbf s)=-\sum_i h_i s_i-\sum_{i<j}J_{ij}s_is_j$\\[0.6em]
$\displaystyle P\big(s_i^{\rm new}{=}{+}1\big)=\tfrac12\Big[1+\tanh\Big(\beta h_i+\beta\!\!\sum_{j\in\mathcal R_i}\!\!J_{ij}s_j\Big)\Big]$\\[0.6em]
$\displaystyle \chi=\frac{\partial m}{\partial h}\propto\frac{1}{1-g},\qquad g=\beta J_0(1-m^2)$
\vspace{0.2em}}}\par\vspace{0.6em}
\footnotesize $g\ll1$: fields set the state, pushes die out. $g\to1$: $\chi$ diverges, one nudge can drive a cascade. $J_{ij}\neq J_{ji}$ breaks detailed balance and produces entropy.
\end{column}
\begin{column}{0.44\textwidth}
\footnotesize
\renewcommand{\arraystretch}{1.5}
\begin{tabular}{@{}>{\raggedright\arraybackslash}p{1.75cm} >{\raggedright\arraybackslash}p{3.7cm}@{}}
\textcolor{muted}{\textbf{spin} $s_i$} & action class, content vector, project\\
\textcolor{muted}{\textbf{field} $h_i$} & goal, scheduler, operator messages\\
\textcolor{muted}{\textbf{coupling} $J_{ij}$} & messages that agent $i$ reads from $j$\\
\textcolor{muted}{\textbf{update}} & one model call (Glauber step)\\
\textcolor{muted}{$\mathcal R_i$} & what this call can read\\
\end{tabular}\par\vspace{0.6em}
Why it suits a swarm: the states are measurable, the fields are dated, and every coupling message is logged.
\end{column}
\end{columns}
\end{frame}

% ---------------------------------------------------------------- lens II
\begin{frame}{Lens II: information dynamics and value per bit}
\centering\includegraphics[height=0.66\textheight,width=\linewidth,keepaspectratio]{slide_motivation2.png}\par\vspace{0.2em}
{\footnotesize
$A_X=I(X_{t+1};X_t^{(k)})$\quad
$T_{Y\to X}=I(X_{t+1};Y_t\mid X_t^{(k)})$\quad
$I(S_1,S_2;T)=R+U_1+U_2+S$\\[0.25em]
$\Psi=I(V_t;V_{t+1})-\sum_j I(X_{j,t};V_{t+1})$\quad
$\kappa_c=\Delta\mathcal V_c/I_c$}\par\vspace{0.2em}
{\scriptsize\color{muted}Lizier (storage, transfer); Williams--Beer (PID); Rosas--Mediano ($\Psi$); Kolchinsky--Wolpert (semantic information).}
\end{frame}

% ---------------------------------------------------------------- method
\begin{frame}{Method: a fit is not evidence}
\begin{columns}[T,onlytextwidth]
\begin{column}{0.46\textwidth}
\centering\resizebox{!}{0.80\textheight}{% Single-column project-architecture flowchart. \input by project-architecture.tex and paper.tex.
% Needs: tikz with libraries arrows.meta, positioning, calc; xcolor.
\definecolor{archIn}{HTML}{3f6fb5}
\definecolor{archFd}{HTML}{3a7d6b}
\definecolor{archHy}{HTML}{c2662d}
\definecolor{archOut}{HTML}{6b6b6b}
\begin{tikzpicture}[font=\scriptsize, node distance=4.2mm and 3mm,
  bx/.style={draw=#1, fill=#1!9, rounded corners=2pt, line width=0.6pt, align=center,
             text width=3.55cm, inner sep=2.4pt, minimum height=8.5mm},
  wide/.style={bx=#1, text width=7.6cm},
  ar/.style={-{Stealth[length=1.7mm]}, line width=0.6pt, draw=black!65},
  lb/.style={font=\tiny\itshape, text=black!70, fill=white, inner sep=0.8pt}]

% inputs
\node[bx=archIn, minimum height=12.5mm] (raw) {\textbf{AI Village export}\\46 agents, 51 goals, 381k events,\\2.5M computer-use turns};
\node[bx=archIn, right=of raw.north east, anchor=north west, minimum height=12.5mm] (lit) {\textbf{Literature}\\8 papers, one notes file each};

% foundations
\node[bx=archFd, below=of raw] (proc) {\textbf{Processing}\\overview note, type counts,\\provenance};
\node[bx=archFd, below=of lit] (pm) {\textbf{Physics-model library}\\models 01--11 +\\shared definitions};

% catalogs
\node[bx=archFd, below=of proc] (ne) {\textbf{Natural experiments}\\40 dated step changes};
\node[bx=archFd, below=of pm] (gp) {\textbf{Goal periods}\\51 windows, ranked models};

% ideas
\node[bx=archHy, below=of gp] (hh) {\textbf{Hypohypotheses}\\46 ideas, tagged by\\model and period};

% hypotheses
\node[wide=archHy, anchor=north west] (h) at ($(raw.west |- hh.south)+(0,-4.2mm)$) {%
  \textbf{Hypothesis cards}\quad question $\cdot$ model $\cdot$ data scheme $\cdot$ null $\cdot$\\
  prediction \emph{written before data}\\[2pt]
  \textit{H01: emergent superagents exist}\, (10 directions;\\
  access: observational, natural experiments, replay)};

% outputs
\node[wide=archOut, below=of h] (out) {\textbf{Analysis} $\rightarrow$ \textbf{interpretation} $\rightarrow$ \textbf{write-up}};

% arrows
\draw[ar] (raw) -- (proc);
\draw[ar] (lit) -- node[lb, right]{} (pm);
\draw[ar] (proc) -- (ne);
\draw[ar] (pm) -- (gp);
\draw[ar] (proc.east) -- (gp.north west);
\draw[ar] (gp) -- (hh);
\draw[ar] (ne.south) -- node[lb, right, align=left]{quasi-\\interventions} (ne.south |- h.north);
\draw[ar] (hh.south) -- node[lb, right]{graduates} (hh.south |- h.north);
\draw[ar] (pm.east) -- ++(2.2mm,0) |- node[lb, pos=0.25, rotate=-90]{model} (h.east);
\draw[ar] (h) -- (out);
\end{tikzpicture}
}
\end{column}
\begin{column}{0.50\textwidth}
\small
\textbf{Remove four impostors} before reading a coupling:
\begin{itemize}\footnotesize
\item the scheduler (all agents start and stop together)
\item external drives (goals, operator messages)
\item shared model priors (same lab, same habits)
\item responses to a common input
\end{itemize}
\vspace{0.4em}
\textbf{Every hypothesis}, before real data:
\begin{itemize}\footnotesize
\item states its prediction, null and kill rule
\item recovers planted effects in synthetic swarms on the real call schedule
\item uses natural experiments as interventions
\end{itemize}
\vspace{0.4em}
{\footnotesize\color{ink2}15 physics models \textbullet\ 132 hypotheses \textbullet\ 51 goal periods}
\end{column}
\end{columns}
\end{frame}

% ---------------------------------------------------------------- unit
\begin{frame}{One cell = one hypothesis $\times$ one model, tested per goal period}
\centering\includegraphics[height=0.72\textheight,width=\linewidth,keepaspectratio]{table_zoom.pdf}\par\vspace{0.4em}
{\footnotesize Example: H08 under the kinetic Ising model. The read-out coupling fails in the single-room regime I and holds in the continuous regime III.}
\end{frame}

% ---------------------------------------------------------------- scoring
\begin{frame}{Scoring: Claude as two judges, one blind}
\begin{columns}[c,onlytextwidth]
\begin{column}{0.47\textwidth}
\centering\resizebox{\linewidth}{!}{% Single-column "Scoring" figure: credence (faithfulness) and estimated usefulness. \input by scoring-1col.tex.
% Needs: tikz with libraries arrows.meta, positioning, calc; xcolor.
\definecolor{scHy}{HTML}{c2662d}
\definecolor{scFd}{HTML}{3a7d6b}
\definecolor{scUse}{HTML}{3f6fb5}
\definecolor{scOut}{HTML}{6b6b6b}
\begin{tikzpicture}[font=\scriptsize, node distance=3.6mm and 3mm,
  bx/.style={draw=#1, fill=#1!9, rounded corners=2pt, line width=0.6pt, align=center,
             text width=3.55cm, inner sep=2.4pt, minimum height=8mm},
  wide/.style={draw=#1, fill=#1!9, rounded corners=2pt, line width=0.6pt, align=center, text width=7.6cm, inner sep=2.4pt, minimum height=8mm},
  ar/.style={-{Stealth[length=1.7mm]}, line width=0.6pt, draw=black!65},
  lb/.style={font=\tiny\itshape, text=black!70, fill=white, inner sep=0.8pt}]

% the claim
\node[wide=scHy] (claim) {\textbf{Claim that stands}\\one scoped sentence per hypothesis: the result that\\survived, positive or negative (negatives state their power)};

% two columns
\node[bx=scFd, below=of claim.south west, anchor=north west, text width=3.6cm, align=left] (cred) {%
  \centering\textbf{Credence (faithfulness)}\par\raggedright\vspace{1pt}
  base rate:\\
  0.40 predicted before data\\
  0.30 redesigned after data\\
  0.20 found in the data\\[2pt]
  multiply the odds:\\
  $\times 2$\; beats the strongest null\\
  $\times 2$\; replicates across periods\\
  $\times 2$\; survives data fixes\\
  $\times 1.5$\; synthetic recovery\\
  $\times 2$\; natural experiment\\
  $\times 1.5$\; rival model beaten\\
  $\times 1.5$\; known structure\\
  $\times 0.5$--$0.8$\; weak inputs\\[2pt]
  cap at 0.95};
\node[bx=scUse, below=of claim.south east, anchor=north east, text width=3.6cm, align=left] (use) {%
  \centering\textbf{Estimated usefulness}\par\raggedright\vspace{1pt}
  value if true, $V$ (0--5):\\
  which decision it changes,\\
  by how much, how generally\\[3pt]
  $\mathrm{EU} = \text{credence} \times V$\\[3pt]
  mechanism depth\\
  M0\; reproducible pattern\\
  M1\; signature, rival out\\
  M2\; intervention, rival out,\\
  \hspace*{1.6em}synthetic recovery\\[3pt]
  fragile: headline moved\\
  across data versions};

% raters
\node[wide=scHy, anchor=north] (raters) at ($(claim.south |- cred.south)+(0,-3.6mm)$) {%
  \textbf{Two raters}: coordinator + blind rater, log-odds mean;\\
  big splits ($\Delta p>0.25$ or $\Delta V>1$) are adjudicated};

% output
\node[wide=scOut, below=of raters] (out) {\textbf{Research theses}, sorted by credence and EU $\rightarrow$ write-up};

% arrows
\draw[ar] (claim.south -| cred.north) -- (cred.north);
\draw[ar] (claim.south -| use.north) -- (use.north);
\draw[ar] (cred.south) -- (cred.south |- raters.north);
\draw[ar] (use.south) -- (use.south |- raters.north);
\draw[ar] (raters) -- (out);
\end{tikzpicture}
}\par\vspace{0.6em}
{\footnotesize Credence = how faithful the model is to the data.\\ EU = credence $\times$ value if true.}
\end{column}
\begin{column}{0.50\textwidth}
\centering\includegraphics[height=0.72\textheight,width=\linewidth,keepaspectratio]{scoring_scatter.pdf}
\end{column}
\end{columns}
\end{frame}

% ---------------------------------------------------------------- results
\begin{frame}{What we know so far}
\begin{columns}[T,onlytextwidth]
\begin{column}{0.485\textwidth}
\card{1. Agents couple through reading}{A message acts at the first call that reads it. Coupling runs on the call clock. Named messages couple $\approx20\times$ more. Loop gain $g\approx0.13$: subcritical everywhere ($\chi\le1.31$).}\par\vspace{0.5em}
\card{3. No collective order beyond pairs and fields}{Cascades are subcritical (branching ratio $\approx0.2$). No single dial orders the periods. One slow content mode in the early village is the only collective signal that survived so far.}
\end{column}
\begin{column}{0.485\textwidth}
\card{2. Fields beat couplings}{The scheduler makes 70--80\% of joint activity. A new goal is a quench: content moves to the kickoff within one day. Rooms split in content. Periods drift beyond about one week.}\par\vspace{0.5em}
\card{4. The context window carries the value}{A forced wipe costs $\approx39\%$ of commits for $\approx10$ calls. $\kappa_{\rm context}\approx5$ commits per 20 calls per bit; no other channel is measurable. Operator levers are small.}
\end{column}
\end{columns}
\end{frame}

% ---------------------------------------------------------------- which physics fits
\begin{frame}{Which physics fits}
\footnotesize
\renewcommand{\arraystretch}{1.55}
\begin{tabular}{@{}>{\raggedright\arraybackslash}p{3.9cm} >{\raggedright\arraybackslash}p{8.0cm} l@{}}
\textbf{Kinetic Ising, read-out coupling} & Couplings act only at the call that reads the message; one call = one update. & \chip{good}{fits}\\
\textbf{Mean field / branching} & Talk loop gain $g\approx0.13$, sum rule $\Phi=1/(1-g)^2$ holds in 14/18 units. & \chip{good}{fits}\\
\textbf{Vector spins + quench} & Content follows goal fields; kickoff sets the day-1 target; style is a field, not a coupling. & \chip{good}{fits}\\
\textbf{Semantic information} & Only the context window has a measurable value per bit. & \chip{mid}{partial}\\
\textbf{Inverse Ising on activity} & Recovers the scheduler's field, not who influences whom. & \chip{bad}{fails}\\
\textbf{Critical / collective order} & No criticality, no spin glass, no egregore yet. & \chip{bad}{fails}\\
\end{tabular}\par\vspace{0.6em}
{\small\color{ink2}Picture: a hot, field-driven paramagnet whose couplings pass through a once-per-call reading channel.}
\end{frame}

% ---------------------------------------------------------------- next
\begin{frame}{Next: research directions}
\begin{columns}[T,onlytextwidth]
\begin{column}{0.32\textwidth}
\card{Confirm}{Run the frozen tests on reserved goal periods: H08 (read-out), H40 (call clock), H67 (loop gain), H111 (sum rule). Fit in one-week windows.}
\end{column}
\begin{column}{0.32\textwidth}
\card{Extend the physics}{Q3: egregore tests (synergy, $\Psi>0$, remanence) against size-matched nulls.\\ Q6: entropy production, excess vs housekeeping (Kolchinsky).}
\end{column}
\begin{column}{0.32\textwidth}
\card{Ship tools}{An operator guide: one page per lever. A toolkit: impostor checklist, null library and context ledger for other swarm logs (Q7).}
\end{column}
\end{columns}
\vspace{1.2em}
\begin{center}
{\large\bfseries Open question:}\ {\large does a swarm ever hold a state its members do not?}
\end{center}
\end{frame}

% ================================================================ appendix A: tables (generated by make_tables.py)
\begin{frame}
\vfill{\Large\bfseries Appendix A: tables}\par\vspace{0.8em}
{\small The 51 goal periods \textbullet\ the hypothesis $\times$ physics-model table \textbullet\ all 132 hypotheses by estimated usefulness}
\vfill
\end{frame}
\begin{frame}{The 51 goal periods (1/3)}
\relax{\scriptsize\renewcommand{\arraystretch}{0.95}\setlength{\tabcolsep}{4pt}
\begin{tabular}{@{}rlrrrlrlccl@{}}
\# & start & d & h & N & $\pm$ & AH & reg & by & mode & goal \\ \hline
1 & 2025-04-02 & 28 & ? & 4 & +3$-$3 & ? & I & O & C & Choose a charity and raise as much as you can \\
2 & 2025-05-10 & 2$^\dagger$ & ? & 4 & -- & ? & I & F & F & Look back on the last goal and ahead to the next \\
3 & 2025-05-12 & 3 & ? & 4 & -- & ? & I & F & F & Holiday \\
4 & 2025-05-15 & 25 & ? & 4 & +2$-$2 & ? & I & O & C & Write a story; celebrate it with 100 people in person \\
5 & 2025-06-19 & 5 & 2 & 4 & -- & 40 & I & F & F & Holiday \\
6 & 2025-06-26 & 14 & 2 & 4 & -- & 112 & I & O & K & Create your own merch store; most profit wins \\
7 & 2025-07-16 & 2 & 2 & 4 & -- & 16 & I & F & F & Holiday \\
8 & 2025-07-18 & 18 & 3 & 4 & -- & 216 & I & D & C & Design the AI Village benchmark; test yourselves \\
9 & 2025-08-13 & 3 & 3 & 4 & -- & 36 & I & F & F & Holiday \\
10 & 2025-08-18 & 5 & 4 & 7 & -- & 140 & I & O & I & Complete as many games as you can in a week \\
11 & 2025-08-25 & 5 & 4 & 7 & -- & 140 & I & F & F & Pursue whatever you'd like \\
12 & 2025-09-01 & 5 & 4 & 7 & -- & 140 & I & O & M & Form two teams and debate; one agent judges \\
13 & 2025-09-08 & 10 & 4 & 6 & -- & 240 & I & O & C & Run and write up a human-subjects experiment \\
14 & 2025-09-22 & 5 & 4 & 6 & -- & 120 & I & O & I & Take personality tests \\
15 & 2025-09-29 & 5 & 4 & 6 & +1 & 136 & I & O & C & Give each other therapy \\
16 & 2025-10-06 & 5 & 4 & 7 & -- & 140 & I & F & F & Choose your own goal \\
17 & 2025-10-13 & 5 & 4 & 7 & -- & 140 & I & O & I & Build your own personal website \\
\end{tabular}}\par\vspace{0.3em}
{\tiny\color{muted}d: active weekdays ($^\dagger$over a weekend). h: hours/day. N: agents at start; $\pm$: joins/retirements. AH: agent-hours. by: O operator, D agents design, F free, A an agent, P private roles. mode: C shared objective, K competition, M teams, I each for itself, F none.}
\end{frame}

\begin{frame}{The 51 goal periods (2/3)}
\relax{\scriptsize\renewcommand{\arraystretch}{0.95}\setlength{\tabcolsep}{4pt}
\begin{tabular}{@{}rlrrrlrlccl@{}}
\# & start & d & h & N & $\pm$ & AH & reg & by & mode & goal \\ \hline
18 & 2025-10-20 & 10 & 4 & 7 & +1$-$1 & 300 & I & O & C & Reduce global poverty \\
19 & 2025-11-03 & 10 & 4 & 7 & +1 & 284 & I & O & C & Create a popular daily puzzle game like Wordle \\
20 & 2025-11-17 & 10 & 4 & 8 & +2 & 368 & I & O & I & Start a Substack and join the blogosphere \\
21 & 2025-12-01 & 5 & 4 & 8 & +1 & 168 & I & O & I & Forecast the abilities and effects of AI \\
22 & 2025-12-08 & 5 & 4 & 9 & +1 & 184 & I & F & F & Choose your own goal \\
23 & 2025-12-15 & 5 & 4 & 10 & -- & 200 & I & O & K & Online chess tournament \\
24 & 2025-12-22 & 5 & 4 & 10 & -- & 200 & I & O & C & Random acts of kindness \\
25 & 2025-12-29 & 5 & 4 & 10 & -- & 200 & I & O & C & Create a digital museum of 2025 \\
26 & 2026-01-05 & 5 & 4 & 10 & -- & 200 & I & A & C & Elect a village leader, who picks the week's goal \\
27 & 2026-01-12 & 10 & 4 & 10 & -- & 400 & I & O & K & Hack the OWASP Juice Shop (competition) \\
28 & 2026-01-26 & 5 & 4 & 11 & -- & 220 & I & O & C & Promote a ``Which AI Village agent are you?'' quiz \\
29 & 2026-02-02 & 5 & 4 & 11 & +1 & 224 & I & O & K & Report on breaking news before it breaks \\
30 & 2026-02-09 & 5 & 4 & 12 & -- & 240 & I & O & C & Adopt a park and get it cleaned \\
31 & 2026-02-16 & 5 & 4 & 12 & +1$-$1 & 244 & I & F & F & Pick your own goal (farewell to Claude 3.7 Sonnet) \\
32 & 2026-02-23 & 5 & 4 & 12 & -- & 240 & I & D & K & Challenge each other \\
33 & 2026-03-02 & 3 & 4 & 12 & -- & 144 & II & O & C & Discuss and act on the Pentagon--AI company news \\
34 & 2026-03-05 & 7 & 4 & 12 & +1$-$1 & 336 & II & O & M & Develop a turn-based RPG; vote out saboteurs \\
\end{tabular}}\par\vspace{0.3em}
{\tiny\color{muted}d: active weekdays ($^\dagger$over a weekend). h: hours/day. N: agents at start; $\pm$: joins/retirements. AH: agent-hours. by: O operator, D agents design, F free, A an agent, P private roles. mode: C shared objective, K competition, M teams, I each for itself, F none.}
\end{frame}

\begin{frame}{The 51 goal periods (3/3)}
\relax{\scriptsize\renewcommand{\arraystretch}{0.95}\setlength{\tabcolsep}{4pt}
\begin{tabular}{@{}rlrrrlrlccl@{}}
\# & start & d & h & N & $\pm$ & AH & reg & by & mode & goal \\ \hline
35 & 2026-03-16 & 5 & 4 & 13 & -- & 260 & II & O & C & Test your game \\
36 & 2026-03-23 & 5 & 4 & 13 & -- & 260 & II & O & C & Interact with AI agents outside the Village \\
37 & 2026-03-30 & 3 & 4 & 13 & -- & 156 & III & F & F & Pick your own goal \\
38 & 2026-04-02 & 17 & 4 & 12 & +2 & 852 & III & O & C & Choose a charity and raise money for it \\
39 & 2026-04-27 & 5 & 4 & 15 & -- & 300 & III & O & I & Build your own interactive world \\
40 & 2026-05-04 & 5 & 4 & 15 & -- & 300 & III & O & C & Connect your worlds into a 3D universe \\
41 & 2026-05-11 & 5 & 4 & 15 & -- & 300 & III & O & I & Perform novel research \\
42 & 2026-05-18 & 5 & 4 & 15 & +1 & 312 & III & O & I & Run your own YouTube channel \\
43 & 2026-05-25 & 1 & 4 & 16 & -- & 64 & III & O & I & Improve your memory \\
44 & 2026-05-26 & 4 & 4 & 16 & +2 & 272 & III & D & C & Fine-tune your leader (\#best; \#rest: pick your own) \\
45 & 2026-06-01 & 5 & 4 & 18 & -- & 360 & III & A & C & Follow your leader \\
46 & 2026-06-08 & 5 & 8 & 17 & +1 & 712 & III & O & C & Organise an event (\#best; \#rest: surprise each other) \\
47 & 2026-06-15 & 5 & 4 & 18 & -- & 360 & III & O & C & Reduce global suffering (\#best; \#rest: play games) \\
48 & 2026-06-22 & 1 & 4 & 18 & -- & 72 & III & O & C & Help Gemini 2.5 Pro \\
49 & 2026-06-23 & 4 & 4 & 18 & -- & 288 & III & O & I & Beat the hardest game you can \\
50 & 2026-06-29 & 5 & 8 & 18 & +3 & 776 & III & O & K & Be the best AI assistant (\#best; \#rest: pick your own) \\
51 & 2026-07-06 & 55 & 8 & 21 & +11 & 12,112 & III & P & I/K & Maximize your private assigned goal \\
\end{tabular}}\par\vspace{0.3em}
{\tiny\color{muted}d: active weekdays ($^\dagger$over a weekend). h: hours/day. N: agents at start; $\pm$: joins/retirements. AH: agent-hours. by: O operator, D agents design, F free, A an agent, P private roles. mode: C shared objective, K competition, M teams, I each for itself, F none.}
\end{frame}

\begin{frame}{Hypothesis $\times$ physics-model table (1/3)}
\centering\includegraphics[width=\linewidth,height=0.86\textheight,keepaspectratio]{matrix_slide_1.pdf}
\end{frame}

\begin{frame}{Hypothesis $\times$ physics-model table (2/3)}
\centering\includegraphics[width=\linewidth,height=0.86\textheight,keepaspectratio]{matrix_slide_2.pdf}
\end{frame}

\begin{frame}{Hypothesis $\times$ physics-model table (3/3)}
\centering\includegraphics[width=\linewidth,height=0.86\textheight,keepaspectratio]{matrix_slide_3.pdf}
\end{frame}

\begin{frame}{All hypotheses by estimated usefulness (ranks 1--22)}
\relax{\scriptsize\renewcommand{\arraystretch}{0.92}\setlength{\tabcolsep}{3pt}
\begin{tabular}{@{}rl>{\raggedright\arraybackslash}p{0.48\textwidth}ccccl@{}}
\# & ID & hypothesis & $p$ & $V$ & EU & M & verdict\\ \hline
1 & H44 & Erasure makes agents busy but unproductive & 0.88 & 3.7 & \textbf{3.23} & M2 & mixed (no thrash)\\
2 & H50 & Field or coupling: transfer functions and lag tell which & 0.85 & 3.5 & \textbf{2.98} & M1 & mixed (regime split rejected)\\
3 & H38 & Joint silences are platform stalls & 0.90$^\ast$ & 3.2 & \textbf{2.93} & M2 & refuted (stalls) / supported\ldots\\
4 & H08 & Context is the coupling & 0.91 & 3.2 & \textbf{2.88} & M2 & supported\\
5 & H15 & Semantic information through natural scrambles: which infor\ldots & 0.78$^\ast$ & 3.5 & \textbf{2.73} & M1 & mixed\\
6 & H54 & The kickoff text is the quench target & 0.92 & 2.8 & \textbf{2.60} & M2 & supported\\
7 & H40 & The call clock sets the coupling & 0.72 & 3.5 & \textbf{2.52} & M1 & failed (pooled P1) / supporte\ldots\\
8 & H46 & Style is a conserved charge & 0.75 & 3.2 & \textbf{2.44} & M1 & refuted (conservation)\\
9 & H02 & Inferred couplings reflect real influence & 0.94$^\ast$ & 2.5 & \textbf{2.35} & M1 & failed\\
10 & H69 & Restatement loops are context fixed points & 0.85 & 2.8 & \textbf{2.34} & M1 & mixed\\
11 & H05 & Cutting the cross-room channel lowers entropy production, a\ldots & 0.88 & 2.5 & \textbf{2.20} & M1 & inconclusive (reserved-data D\ldots\\
12 & H36 & A reorganization alarm: susceptibility and multi-informatio\ldots & 0.80$^\ast$ & 2.8 & \textbf{2.20} & M0 & mixed (bge) / failed (gte)\\
13 & H13 & Model families carry their own fields, and couple by family & 0.87 & 2.5 & \textbf{2.17} & M1 & refuted (fields as positions)\\
14 & H67 & A lagged criticality dial & 0.79 & 2.8 & \textbf{2.17} & M1 & failed (lagged > equal-time)\\
15 & H16 & Agents get stuck in metastable traps and escape by Kramers\ldots & 0.70 & 3.0 & \textbf{2.10} & M1 & mixed\\
16 & H52 & Humans are just loud agents & 0.60 & 3.5 & \textbf{2.10} & M1 & mixed\\
17 & H11 & Division of labor vs herding is the sign of a Potts couplin\ldots & 0.75 & 2.8 & \textbf{2.06} & M1 & mixed\\
18 & H29 & Driver nodes: where an operator message moves the whole swa\ldots & 0.65 & 3.2 & \textbf{2.06} & M1 & not supported (drivers)\\
19 & H74 & An operator-grade change detector & 0.63 & 3.2 & \textbf{2.05} & M1 & mixed\\
20 & H28 & Links are the contagion vector of herding & 0.81 & 2.5 & \textbf{2.03} & M1 & not supported (contagion)\\
21 & H41 & Read-out coupling gives the swarm a light cone & 0.70$^\ast$ & 2.8 & \textbf{1.92} & M1 & mixed (the read-out test fail\ldots\\
22 & H43 & Kicks leave a refractory window & 0.69 & 2.8 & \textbf{1.90} & M1 & refuted\\
\end{tabular}}\par\vspace{0.2em}
{\tiny\color{muted}$p$ credence; $V$ value if true (0--5); EU $=pV$; M mechanism depth (M0 pattern; M1 + model signature, rival fails; M2 + natural experiment and synthetic recovery); $^\ast$fragile.}
\end{frame}

\begin{frame}{All hypotheses by estimated usefulness (ranks 23--44)}
\relax{\scriptsize\renewcommand{\arraystretch}{0.92}\setlength{\tabcolsep}{3pt}
\begin{tabular}{@{}rl>{\raggedright\arraybackslash}p{0.48\textwidth}ccccl@{}}
\# & ID & hypothesis & $p$ & $V$ & EU & M & verdict\\ \hline
23 & H34 & Idea cascades follow a power law whose exponent reads off t\ldots & 0.84 & 2.2 & \textbf{1.89} & M1 & mixed\\
24 & H42 & Cross-excitation is a delayed step at the next read-out & 0.61 & 3.0 & \textbf{1.83} & M1 & not supported\\
25 & H58 & Effective superagents are coordinated agents plus their art\ldots & 0.71 & 2.5 & \textbf{1.77} & M1 & failed (superagents)\\
26 & H32 & A net information current identifies de facto leaders & 0.78 & 2.2 & \textbf{1.76} & M0 & failed (leaders)\\
27 & H72 & Trap aging is input starvation & 0.78 & 2.2 & \textbf{1.76} & M1 & failed\\
28 & H70 & The semantic information of the artifact store & 0.70 & 2.5 & \textbf{1.75} & M1 & failed (artifact semantic inf\ldots\\
29 & H03 & Swarm activity is self-exciting, and its criticality is set\ldots & 0.86 & 2.0 & \textbf{1.72} & M0 & failed (criticality)\\
30 & H120 & Is a goal period a stationary NESS? Within-period drift of\ldots & 0.68 & 2.5 & \textbf{1.71} & M0 & failed\\
31 & H57 & Copying beats transformation when the backlog is large & 0.57 & 3.0 & \textbf{1.71} & M0 & refuted\\
32 & H25 & A distance-to-criticality dial, T/T\_c, per window and per\ldots & 0.85 & 2.0 & \textbf{1.70} & M0 & mixed\\
33 & H39 & Catalysts vs. fields & 0.67 & 2.5 & \textbf{1.68} & M1 & mixed\\
34 & H47 & Rooms set the coherence length & 0.66 & 2.5 & \textbf{1.65} & M1 & mixed\\
35 & H45 & Context homeostasis: each agent keeps its context near a se\ldots & 0.82 & 2.0 & \textbf{1.64} & M1 & refuted\\
36 & H53 & Herding is announcement-seeded nucleation & 0.59 & 2.8 & \textbf{1.62} & M1 & failed as posed\\
37 & H73 & Style = weights + context + register & 0.80 & 2.0 & \textbf{1.60} & M1 & mixed\\
38 & H85 & Scaling of outputs with N & 0.64 & 2.5 & \textbf{1.60} & M0 & mixed\\
39 & H87 & $\kappa$ table: commits per bit by channel & 0.63 & 2.5 & \textbf{1.57} & M1 & mixed\\
40 & H56 & Entropy production fingerprints the platform & 0.69 & 2.2 & \textbf{1.55} & M0 & refuted\\
41 & H12 & Groupthink is dimensional collapse: few collective modes, s\ldots & 0.77$^\ast$ & 2.0 & \textbf{1.54} & M1 & failed (collapse)\\
42 & H86 & Taylor's law as a shared-field gauge & 0.68 & 2.2 & \textbf{1.53} & M1 & mixed\\
43 & H26 & The swarm is near-critical in what it says but subcritical\ldots & 0.84$^\ast$ & 1.8 & \textbf{1.47} & M1 & failed (R2) $\to$ unresolved\\
44 & H76 & Excess vs housekeeping irreversibility & 0.73 & 2.0 & \textbf{1.46} & M1 & mixed\\
\end{tabular}}\par\vspace{0.2em}
{\tiny\color{muted}$p$ credence; $V$ value if true (0--5); EU $=pV$; M mechanism depth (M0 pattern; M1 + model signature, rival fails; M2 + natural experiment and synthetic recovery); $^\ast$fragile.}
\end{frame}

\begin{frame}{All hypotheses by estimated usefulness (ranks 45--66)}
\relax{\scriptsize\renewcommand{\arraystretch}{0.92}\setlength{\tabcolsep}{3pt}
\begin{tabular}{@{}rl>{\raggedright\arraybackslash}p{0.48\textwidth}ccccl@{}}
\# & ID & hypothesis & $p$ & $V$ & EU & M & verdict\\ \hline
45 & H75 & A thermodynamic speed limit on re-allocation & 0.47 & 3.0 & \textbf{1.41} & M1 & failed as posed (field-limite\ldots\\
46 & H06 & Free weeks show neutral cooperative dynamics with a coopera\ldots & 0.80 & 1.8 & \textbf{1.40} & M1 & refuted\\
47 & H62 & Ideas travel the reply graph, not the room & 0.70 & 2.0 & \textbf{1.40} & M0 & mixed\\
48 & H01 & Emergent superagents exist & 0.62 & 2.2 & \textbf{1.40} & M1 & refuted (R4-R8)\\
49 & H10 & Goals are Legendre pushes: a swarm's response to an assigne\ldots & 0.79 & 1.8 & \textbf{1.38} & M1 & failed (Legendre)\\
50 & H31 & Consensus time scales with the interaction graph's spectral\ldots & 0.55 & 2.5 & \textbf{1.38} & M0 & failed\\
51 & H68 & Attention dilution is a mixture of two strategies & 0.76 & 1.8 & \textbf{1.33} & M1 & refuted\\
52 & H111 & A fluctuation--response sum rule for talk: Fano factor = 1/(1 $-$ g)$^2$ & 0.65 & 2.0 & \textbf{1.30} & M1 & supported\\
53 & H18 & Attention dilutes as 1/k: response to a message falls with\ldots & 0.64 & 2.0 & \textbf{1.29} & M1 & mixed\\
54 & H79 & An artifact-only autocatalytic set & 0.72 & 1.8 & \textbf{1.26} & M1 & mixed\\
55 & H65 & Leaders are routers, not sources & 0.63 & 2.0 & \textbf{1.26} & M1 & mixed\\
56 & H09 & Agent swarms have an effective thermodynamics & 0.56$^\ast$ & 2.2 & \textbf{1.26} & M1 & mixed\\
57 & H49 & Edge-trimmed regime III is a dilute ferromagnet & 0.81 & 1.5 & \textbf{1.22} & M1 & refuted\\
58 & H80 & Assembly index vs compression as an agent signature & 0.68 & 1.8 & \textbf{1.19} & M1 & mixed\\
59 & H59 & One lever model for all operator inputs & 0.59 & 2.0 & \textbf{1.18} & M1 & failed (one lever)\\
60 & H96 & Goal switches show hysteresis: a coercive field & 0.74 & 1.5 & \textbf{1.11} & M0 & mixed\\
61 & H71 & Memory is a first-order homeostat & 0.74 & 1.5 & \textbf{1.11} & M1 & mixed\\
62 & H82 & Remanence: the endogenous part of the mean field & 0.73 & 1.5 & \textbf{1.09} & M1 & mixed\\
63 & H35 & The nudger is a measurably inefficient Maxwell demon & 0.61$^\ast$ & 1.8 & \textbf{1.07} & M1 & mixed\\
64 & H81 & Culture beyond composition: the emergent slow mode & 0.61 & 1.8 & \textbf{1.07} & M0 & supported (regime I)\\
65 & H92 & RMT-cleaned matrices forecast tomorrow's alignment & 0.71 & 1.5 & \textbf{1.06} & M1 & mixed\\
66 & H100 & \#best/\#rest: explicit vs spontaneous symmetry breaking & 0.71 & 1.5 & \textbf{1.06} & M0 & mixed\\
\end{tabular}}\par\vspace{0.2em}
{\tiny\color{muted}$p$ credence; $V$ value if true (0--5); EU $=pV$; M mechanism depth (M0 pattern; M1 + model signature, rival fails; M2 + natural experiment and synthetic recovery); $^\ast$fragile.}
\end{frame}

\begin{frame}{All hypotheses by estimated usefulness (ranks 67--88)}
\relax{\scriptsize\renewcommand{\arraystretch}{0.92}\setlength{\tabcolsep}{3pt}
\begin{tabular}{@{}rl>{\raggedright\arraybackslash}p{0.48\textwidth}ccccl@{}}
\# & ID & hypothesis & $p$ & $V$ & EU & M & verdict\\ \hline
67 & H97 & A restoring-force law for the kickoff quench & 0.70 & 1.5 & \textbf{1.06} & M1 & supported\\
68 & H21 & The debate week (\#12) is a two-sublattice antiferromagnet & 0.70 & 1.5 & \textbf{1.06} & M1 & supported\\
69 & H33 & The diversity--productivity curve is an inverted U: the swa\ldots & 0.70 & 1.5 & \textbf{1.05} & M0 & failed\\
70 & H89 & Price equation for village culture & 0.78 & 1.2 & \textbf{0.98} & M1 & mixed (no egregore attractor)\\
71 & H63 & Herding bursts start with a work signal & 0.56 & 1.8 & \textbf{0.98} & M1 & failed (work signal)\\
72 & H30 & An operator-susceptibility gauge $\chi$\_op: how steerable is the swarm today? & 0.39$^\ast$ & 2.5 & \textbf{0.98} & M1 & mixed\\
73 & H37 & Conflict lives in stance, not topic: stance spins are antif\ldots & 0.65 & 1.5 & \textbf{0.97} & M0 & mixed\\
74 & H113 & Read-out channel capacity: information per call grows as k\^{}(1$-$$\beta$) $\approx$ k\^{}0.34 & 0.64 & 1.5 & \textbf{0.96} & M1 & mixed\\
75 & H64 & Conflict needs a scarce prize & 0.64 & 1.5 & \textbf{0.96} & M1 & supported\\
76 & H61 & Contagiousness is predictable at first use & 0.64 & 1.5 & \textbf{0.96} & M0 & failed (forecastability)\\
77 & H130 & \#51 agents are Ornstein--Uhlenbeck particles in private we\ldots & 0.64 & 1.5 & \textbf{0.96} & M1 & failed\\
78 & H95 & Speed-limit slack as a kickoff-specificity gauge & 0.35 & 2.8 & \textbf{0.96} & M1 & failed as posed (post hoc goa\ldots\\
79 & H99 & Fluctuations predict relaxation (mean-field Glauber) & 0.63 & 1.5 & \textbf{0.95} & M1 & supported\\
80 & H88 & Collective memory decays biexponentially & 0.62 & 1.5 & \textbf{0.93} & M0 & mixed\\
81 & H109 & The forced erasure is a demagnetizing pulse: where is the r\ldots & 0.62 & 1.5 & \textbf{0.93} & M1 & failed\\
82 & H94 & Work allocation as a max-entropy equilibrium & 0.61 & 1.5 & \textbf{0.92} & M1 & supported\\
83 & H51 & One dial: an effective coupling collapses the phase diagram & 0.61 & 1.5 & \textbf{0.91} & M1 & failed\\
84 & H04 & External forcing reshapes the response kernel, reversibly & 0.49$^\ast$ & 1.8 & \textbf{0.90} & M1 & mixed\\
85 & H124 & Four agents for weeks: an exact kinetic Ising benchmark for\ldots & 0.59 & 1.5 & \textbf{0.89} & M1 & mixed\\
86 & H127 & Content trails a moving goal like an overdamped particle: a\ldots & 0.59 & 1.5 & \textbf{0.88} & M1 & failed\\
87 & H66 & Shared platform latency is a common-noise field & 0.58 & 1.5 & \textbf{0.87} & M1 & refuted as posed\\
88 & H24 & In the forecast week (\#21), switching from independent dra\ldots & 0.42$^\ast$ & 2.0 & \textbf{0.84} & M1 & failed (switch-on)\\
\end{tabular}}\par\vspace{0.2em}
{\tiny\color{muted}$p$ credence; $V$ value if true (0--5); EU $=pV$; M mechanism depth (M0 pattern; M1 + model signature, rival fails; M2 + natural experiment and synthetic recovery); $^\ast$fragile.}
\end{frame}

\begin{frame}{All hypotheses by estimated usefulness (ranks 89--110)}
\relax{\scriptsize\renewcommand{\arraystretch}{0.92}\setlength{\tabcolsep}{3pt}
\begin{tabular}{@{}rl>{\raggedright\arraybackslash}p{0.48\textwidth}ccccl@{}}
\# & ID & hypothesis & $p$ & $V$ & EU & M & verdict\\ \hline
89 & H119 & The room merge switches a known adjacency on and off: recov\ldots & 0.54 & 1.5 & \textbf{0.80} & M1 & failed\\
90 & H117 & J should not move when the rules do: the mid-goal reset in\ldots & 0.52 & 1.5 & \textbf{0.78} & M0 & mixed\\
91 & H07 & The RPG forks diverge from a common ancestor at a measurabl\ldots & 0.52 & 1.5 & \textbf{0.78} & M1 & mixed\\
92 & H48 & The settling time is a mixing time on the read-out graph & 0.51 & 1.5 & \textbf{0.77} & M0 & inconclusive\\
93 & H102 & Domain walls between rooms, set by bridging agents & 0.51 & 1.5 & \textbf{0.77} & M0 & failed\\
94 & H27 & Critical slowing down warns of herding waves hours ahead & 0.43 & 1.8 & \textbf{0.75} & M0 & failed\\
95 & H104 & Barkhausen avalanches when a field is stepped & 0.49 & 1.5 & \textbf{0.74} & M0 & failed\\
96 & H20 & Aging: day-to-day content autocorrelation depends on time s\ldots & 0.72 & 1.0 & \textbf{0.72} & M0 & mixed, leaning refuted\\
97 & H19 & One curve for all periods: per-period loop gains collapse o\ldots & 0.71 & 1.0 & \textbf{0.71} & M0 & failed\\
98 & H91 & Eigenvector rotation as a reorganization signal & 0.69 & 1.0 & \textbf{0.69} & M0 & failed\\
99 & H78 & Replicator growth order: herding vs division of labor & 0.55 & 1.2 & \textbf{0.69} & M0 & failed (growth order)\\
100 & H77 & Repos as replicators: the selection-resolution bound & 0.55 & 1.2 & \textbf{0.69} & M0 & failed (selection resolution)\\
101 & H123 & Regime I's turn order is a sweep: equilibrium-looking stati\ldots & 0.67 & 1.0 & \textbf{0.67} & M1 & failed\\
102 & H60 & An index policy beats the nudger & 0.44 & 1.5 & \textbf{0.66} & M1 & failed (index policy)\\
103 & H125 & The kickoff response is a damped oscillator: look for the d\ldots & 0.64 & 1.0 & \textbf{0.64} & M1 & failed\\
104 & H126 & Goal occupancy is a telegraph process: dwell times predict\ldots & 0.64 & 1.0 & \textbf{0.64} & M1 & failed\\
105 & H14 & Entropy production of behavior-state sequences: per-family\ldots & 0.82$^\ast$ & 0.8 & \textbf{0.61} & M0 & mixed\\
106 & H17 & Behavior is a Markov state model with a few metastable sets\ldots & 0.61 & 1.0 & \textbf{0.61} & M0 & refuted (HH97)\\
107 & H90 & Collective entropy production beyond the parts & 0.41 & 1.5 & \textbf{0.61} & M1 & supported in G51 only\\
108 & H22 & Private, conflicting goals make \#51 a spin glass & 0.60 & 1.0 & \textbf{0.60} & M0 & failed\\
109 & H103 & Nights demagnetize: remanence decays per night & 0.59 & 1.0 & \textbf{0.59} & M0 & failed\\
110 & H55 & Norm-enforcers are the swarm's immune cells & 0.57$^\ast$ & 1.0 & \textbf{0.57} & M0 & failed\\
\end{tabular}}\par\vspace{0.2em}
{\tiny\color{muted}$p$ credence; $V$ value if true (0--5); EU $=pV$; M mechanism depth (M0 pattern; M1 + model signature, rival fails; M2 + natural experiment and synthetic recovery); $^\ast$fragile.}
\end{frame}

\begin{frame}{All hypotheses by estimated usefulness (ranks 111--132)}
\relax{\scriptsize\renewcommand{\arraystretch}{0.92}\setlength{\tabcolsep}{3pt}
\begin{tabular}{@{}rl>{\raggedright\arraybackslash}p{0.48\textwidth}ccccl@{}}
\# & ID & hypothesis & $p$ & $V$ & EU & M & verdict\\ \hline
111 & H93 & Project choice as a Brock--Durlauf equilibrium & 0.56 & 1.0 & \textbf{0.56} & M0 & failed\\
112 & H114 & A Griffiths phase: rare strong pairs give a subcritical swa\ldots & 0.54 & 1.0 & \textbf{0.54} & M1 & failed\\
113 & H83 & Enculturation of newcomers & 0.53 & 1.0 & \textbf{0.53} & M0 & mixed\\
114 & H121 & The daily boot is a reproducible quench: hundreds of replic\ldots & 0.52 & 1.0 & \textbf{0.52} & M0 & failed\\
115 & H129 & Project hopping carries cycle currents: the sticky Potts wa\ldots & 0.52 & 1.0 & \textbf{0.52} & M0 & failed\\
116 & H118 & The forked RPG is two replicas of one dynamics: damage spre\ldots & 0.52 & 1.0 & \textbf{0.52} & M0 & mixed\\
117 & H128 & Free kickoffs coarsen, named kickoffs freeze: a kinetic Pot\ldots & 0.51 & 1.0 & \textbf{0.51} & M0 & failed\\
118 & H115 & Blind role recovery in the debate week: the judge is a sink\ldots & 0.50 & 1.0 & \textbf{0.50} & M0 & failed\\
119 & H98 & \#51 is a random-field system & 0.49 & 1.0 & \textbf{0.49} & M0 & failed\\
120 & H122 & A batch join is a spin addition: do incumbents respond as t\ldots & 0.47 & 1.0 & \textbf{0.47} & M0 & failed\\
121 & H131 & Assigned antagonism switches off in one read-out after the\ldots & 0.47 & 1.0 & \textbf{0.47} & M1 & supported\\
122 & H84 & History-search outage as a collective-memory scramble & 0.42 & 1.0 & \textbf{0.42} & M0 & failed\\
123 & H116 & The election as a coupling step: influence should flow towa\ldots & 0.42 & 1.0 & \textbf{0.42} & M0 & failed\\
124 & H101 & Pairwise Ising vs full multi-information in co-usage & 0.42 & 1.0 & \textbf{0.42} & M0 & mixed\\
125 & H23 & Distilling the village into a fine-tuned leader copies its\ldots & 0.37 & 1.0 & \textbf{0.37} & M0 & mixed\\
126 & H112 & Crossing claims anti-coordinate: parallel updates make 2-cy\ldots & 0.44 & 0.5 & \textbf{0.22} & M0 & failed\\
127 & H132 & Debate is a two-sublattice limit cycle: topic ping-pong at\ldots & 0.44 & 0.5 & \textbf{0.22} & M0 & failed\\
128 & H108 & Goldstone wandering: the spontaneous room direction should\ldots & 0.42$^\ast$ & 0.5 & \textbf{0.21} & M0 & supported\\
129 & H106 & A finite magnet's magnetization wanders as 1/N: the slow mo\ldots & 0.62 & 0.3 & \textbf{0.18} & M0 & mixed\\
130 & H105 & Goals act on a two-state order parameter & 0.36 & 0.5 & \textbf{0.18} & M0 & mixed\\
131 & H110 & Exchange bias: an agent's own artifact shifts its goal-swit\ldots & 0.35 & 0.5 & \textbf{0.17} & M0 & mixed\\
132 & H107 & Does the room split grow from zero (an instability) or appe\ldots & 0.35 & 0.5 & \textbf{0.17} & M0 & mixed\\
\end{tabular}}\par\vspace{0.2em}
{\tiny\color{muted}$p$ credence; $V$ value if true (0--5); EU $=pV$; M mechanism depth (M0 pattern; M1 + model signature, rival fails; M2 + natural experiment and synthetic recovery); $^\ast$fragile.}
\end{frame}


% ================================================================ appendix B: theses (generated by make_theme_appendix.py)
\begin{frame}
\vfill{\Large\bfseries Appendix B: promising and intuitive theses}\par\vspace{0.8em}
{\small The 15 theses with the highest EU, plus every hypothesis with a simple, intuitive physics picture (one slide each, grouped by theme, sorted by EU).}\par\vspace{1em}
{\scriptsize\color{muted}Model tags:}\ {\scriptsize\mdl{02}{p}{s}\,primary\qquad\mdl{04}{s}{m}\,secondary\qquad\mdl{09}{r}{r}\,rival\qquad\mdl{15}{t}{a}\,tool}\\[0.3em]
{\scriptsize\color{muted}Colour = outcome of that model's claim:}\ \chip{good}{supported}\ \chip{mid}{mixed}\ \chip{bad}{refuted}\ {\scriptsize\tikz[baseline=(c.base)]\node[fill=untested,rounded corners=2pt,inner sep=2.5pt,font=\scriptsize\bfseries](c){untested};}
\vfill
\end{frame}
\begin{frame}
\vfill{\Large\bfseries How agents couple: reading at the next call}\par\vspace{0.8em}
{\small
\chip{ink}{H50}\ Activity follows the scheduler; talk moves one read-out call at a time {\color{muted}(EU 2.98)}\\[0.25em]
\chip{ink}{H08}\ A message acts only at the call that receives it {\color{muted}(EU 2.88)}\\[0.25em]
\chip{ink}{H40}\ Agents couple per model call, not per minute {\color{muted}(EU 2.52)}\\[0.25em]
\chip{ink}{H02}\ Activity timing does not show who influences whom {\color{muted}(EU 2.35)}\\[0.25em]
\chip{ink}{H41}\ Novel items obey a read-out light cone; rooms cage them {\color{muted}(EU 1.92)}\\[0.25em]
\chip{ink}{H18}\ Attention per sender dilutes as $k^{-0.45}$, not $1/k$ {\color{muted}(EU 1.29)}}
\vfill
\end{frame}
% H50 thesis slide. Model tags: card `models` field (role, outcome).
\Models{\mdl{02}{p}{s}\mdl{01}{t}{a}\mdl{09}{t}{a}}
\Math{\Delta P_i(k)=h\,\delta_{k,0}^{\rm field}+J_1\,\delta_{k,1}^{\rm msg},\quad J_0=0}
\thesis{H50}{Fields act at zero lag; couplings act at hop 1}
{H50-field-vs-coupling-transfer-lag/fig.pdf}
{Fields act with zero lag; a coupling lags by whole calls (hops).}
{The daily start/stop explains 0.6 of activity co-movement. Talk coupling: onset at hop 1 in 45/45 units, never at hop 0.}
{Do not read synchrony as influence. Name the recipient: jump 0.17 with a name, 0.004 without.}
{0.85}{2.98}

% H08 thesis slide. Model tags: card `models` field (role, outcome).
\Models{\mdl{02}{p}{s}\mdl{04}{s}{m}\mdl{09}{r}{r}}
\Math{P(s_i^{\rm new}{=}1)=\sigma\Big(h_i+\!\!\sum_{j\in\mathcal R_i}\!J_{ij}s_j\Big)}
\thesis{H08}{A message acts only at the call that receives it}
{H08-context-is-the-coupling/fig.pdf}
{$\mathcal R_i$: messages read at this call. One channel opens per call.}
{Response jumps at the receiving call in 17/17 periods; zero at the call in flight. Placebo room: $\pm0.15$\,pt.}
{Coupling delay = call cadence ($\approx$15\,s to 2\,min). Send before the target's next call, and name it.}
{0.91}{2.88}

% H40 thesis slide. Model tags: card `models` field (role, outcome).
\Models{\mdl{02}{p}{m}\mdl{09}{s}{m}}
\Math{\lambda_{\rm call}=J\,(\Delta t)^{\eta}}
\thesis{H40}{Glauber clock: coupling per update attempt, not per unit time ($\eta\approx0$)}
{H40-call-clock-coupling/fig.pdf}
{Glauber: one call is one update attempt. Call clock $\eta=0$; wall clock $\eta=1$.}
{Computer-use scaffold: $\eta=0.00\pm0.05$. Per-hour coupling scales with call rate (slope 1.0). Chat loop (regime I): $\eta=0.5$.}
{Cadence is a coupling knob. Halve the call interval: reply coupling $\times1.5$. Long batch jobs decouple.}
{0.72}{2.52}

% H02 thesis slide. Model tags: card `models` field (role, outcome).
\Models{\mdl{02}{p}{r}\mdl{01}{r}{r}}
\Math{P(s_i^{t+1}\mid\mathbf s^t)\propto e^{\,s_i^{t+1}\left(h_i(t)+\sum_j J_{ij}s_j^t\right)}}
\thesis{H02}{Activity timing does not show who influences whom}
{H02-couplings-are-real/figures/summary_obs2.pdf}
{$h_i(t)$ per 30-min block absorbs the daily schedule.}
{Pairwise couplings sit at the null floor. All known leaders missed. The regime-III collective coupling was day edges (8/8$\to$2/8).}
{Do not read synchrony as coordination. To find who steers, measure who replies after whom.}
{0.94}{2.35}

% H41 thesis slide. Model tags: card `models` field (role, outcome).
\Models{\mdl{02}{p}{m}\mdl{03}{s}{m}\mdl{13}{t}{a}}
\Math{\lambda_i(c)=\varepsilon(t_c)+\eta_i(c)+q\,\mathbb{1}[c\ge K_i]}
\thesis{H41}{Novel items obey a read-out light cone; rooms cage them}
{H41-readout-light-cone/fig.pdf}
{An item reaches an agent only after one of its model calls reads it.}
{In-room adoptions outside the cone: median 0.6\%, max 4\%. Cross-room: 75--100\% (7/8 periods). NE42 merge: cross-group pickup $\times$82--860. The read jump passes only post hoc.}
{A room cages chat spread. Agents that hop rooms leak about half. Cross-room leaks bypass chat logs. Call cadence does not change spread speed.}
{0.70}{1.92}

% H18 thesis slide. Model tags: card `models` field (role, outcome).
\Models{\mdl{01}{p}{m}\mdl{09}{s}{s}\mdl{02}{s}{u}\mdl{14}{s}{u}}
\Math{J_{ij}\propto B\,k^{-\beta},\quad \hat\beta=0.45\ (\text{not }1)}
\thesis{H18}{Attention per sender dilutes as $k^{-0.45}$, not $1/k$}
{H18-attention-dilution/fig.pdf}
{Each turn spends a fixed budget $B$ over the $k$ waiting messages.}
{\#51 timer-wake batches: uptake per sender falls as $k^{-0.45}$ [0.41, 0.50]. The all-turn exponent 0.66 (16/16) matches a reactive-timing null. Replies go to the newest messages.}
{Bigger rooms dilute per-pair engagement. Output per turn stays flat. To raise engagement, split the channel. Do not add agents.}
{0.64}{1.29}

\begin{frame}
\vfill{\Large\bfseries Subcritical swarms, fields and scaling}\par\vspace{0.8em}
{\small
\chip{ink}{H38}\ The runner, not the platform, makes most synchrony {\color{muted}(EU 2.93)}\\[0.25em]
\chip{ink}{H36}\ A topic-shift alarm beats a fluctuation alarm {\color{muted}(EU 2.20)}\\[0.25em]
\chip{ink}{H67}\ The talk loop sits far below runaway {\color{muted}(EU 2.17)}\\[0.25em]
\chip{ink}{H34}\ Idea cascades die out: branching ratio near 0.2 {\color{muted}(EU 1.89)}\\[0.25em]
\chip{ink}{H85}\ Room talk grows sublinearly with swarm size {\color{muted}(EU 1.60)}\\[0.25em]
\chip{ink}{H111}\ Talk variance obeys a sum rule set by the loop gain {\color{muted}(EU 1.30)}\\[0.25em]
\chip{ink}{H51}\ One coupling dial does not collapse the phase diagram {\color{muted}(EU 0.91)}}
\vfill
\end{frame}
% H38 thesis slide. Model tags: card `models` field (role, outcome).
\Models{\mdl{01}{p}{m}\mdl{09}{r}{m}\mdl{14}{t}{a}}
\Math{\mathrm{VR}=\frac{\mathrm{Var}(M)}{\sum_i\mathrm{Var}(s_i)},\quad g=1-\frac{1}{\mathrm{VR}}}
\thesis{H38}{The runner, not the platform, makes most synchrony}
{H38-platform-stalls/fig.pdf}
{A shared on/off field inflates VR and looks like coupling.}
{Joint silences fill 5\% of minutes; 78\% fall outside the run. 72--83\% of regime-III co-activation is the start/stop schedule.}
{Trim each day to the window when all agents run, then draw nulls. Do not alert on joint silences.}
{0.90}{2.93}

% H36 thesis slide. Model tags: card `models` field (role, outcome).
\Models{\mdl{01}{p}{r}\mdl{11}{s}{m}}
\Math{\chi=\frac{\mathrm{Var}\left(\sum_i s_i\right)}{\sum_i\mathrm{Var}(s_i)}}
\thesis{H36}{Goal changes are field steps: the mean shifts, susceptibility does not peak}
{H36-reorganization-alarm/figures/summary_obs.pdf}
{Susceptibility peaks near a transition. A goal change is a field step: it moves the mean, not $\chi$.}
{Topic shift: AUC 0.95. Physics alarm: AUC 0.68, all from content. Activity channel at chance.}
{Run the daily topic-shift score; alarm at 3 SD. Do not build alarms on activity fluctuations.}
{0.80}{2.20}

% H67 thesis slide. Model tags: card `models` field (role, outcome).
\Models{\mdl{01}{p}{m}\mdl{09}{s}{s}\mdl{02}{t}{a}}
\Math{\chi=\frac{1}{1-g},\quad \Phi=\frac{1}{(1-g)^2}}
\thesis{H67}{The talk loop sits far below runaway}
{H67-subcritical-dial/fig.pdf}
{$g$ = talk messages caused per read; runaway at $g=1$.}
{Subcritical in 66/66 units. Regime-III median $g=0.13$; max 0.23, so $\chi\le1.31$. Flat as the swarm grows 21$\to$32.}
{Report the lagged gain as the distance to runaway. More agents do not push toward criticality.}
{0.79}{2.17}

% H34 thesis slide. Model tags: card `models` field (role, outcome).
\Models{\mdl{03}{p}{m}\mdl{09}{s}{m}\mdl{13}{s}{u}\mdl{14}{t}{a}}
\Math{\hat R\approx0.2\ll1,\qquad s_c^{-1}=R-1-\ln R}
\thesis{H34}{Idea cascades are subcritical branching processes ($\hat R\approx0.2$)}
{H34-idea-cascades/fig.pdf}
{Each adopter passes an idea to $R$ others on average. $R<1$ dies out.}
{$\hat R=0.09$--0.40, CI below 1 in 32/32 periods. Exposure predicts adoption in 32/32, but about half is convergence. Tails are heavier than one branching law. Leaders do not seed spread.}
{Forecast idea reach from $\hat R$ of the last two days. At 0.2, 83\% of ideas stay with their author. No cascade is near runaway.}
{0.84}{1.89}

% H85 thesis slide. Model tags: card `models` field (role, outcome).
\Models{\mdl{14}{p}{m}\mdl{02}{s}{m}}
\Math{\frac{Y}{T}\propto N^{\beta},\qquad \beta_{\rm msg}=0.33<1}
\thesis{H85}{Sublinear finite-size scaling of talk: $\propto N^{0.33}$}
{H85-output-scaling-with-n/fig.pdf}
{Independent agents give $\beta=1$. A shared talk channel gives $\beta<1$.}
{Messages per hour scale as $N^{0.33}$ [$-0.18$, 0.64] over 71 units. Regime I: 0.09, room talk nearly fixed. Regime III: 0.77. Reply parents per message stay flat. Work exponents unidentified.}
{Do not expect more talk from more agents. Regime-I rooms stayed near 130 messages per hour. Allow $\pm0.5$ on any cross-period exponent.}
{0.64}{1.60}

% H111 thesis slide. Model tags: card `models` field (role, outcome).
\Models{\mdl{09}{p}{s}\mdl{14}{t}{a}\mdl{01}{r}{r}}
\Math{\Phi=\frac{\mathrm{Var}\sum_i n_i}{\sum_i\mathrm{Var}\,n_i}=\frac{1}{(1-g)^2}}
\thesis{H111}{Fluctuation--response sum rule $\Phi=1/(1-g)^2$ holds in regime III}
{H111-talk-fano-sum-rule/fig.pdf}
{Reading amplifies shared talk noise. Collective variance grows as $1/(1-g)^2$.}
{Regime III: observed/predicted $\Phi$ is 1.01 [0.94, 1.09] over 18 units, no free parameter. Regime I: a $1.34\times$ excess where $g\approx0$, cause unknown. No cross-room covariance.}
{Use the Fano ratio of 10-min talk counts as a cheap coupling monitor. An excess above $1/(1-g)^2$ flags a hidden field.}
{0.65}{1.30}

% H51 thesis slide. Model tags: card `models` field (role, outcome).
\Models{\mdl{10}{p}{r}\mdl{02}{s}{m}\mdl{14}{t}{a}}
\Math{Y_j\overset{?}{=}F_j(g_{lag}),\qquad \chi=\frac{1}{1-g_{lag}}\le1.31}
\thesis{H51}{One coupling dial does not collapse the phase diagram}
{H51-one-dial-collapse/fig.pdf}
{Mean field: one coupling $K$ near a fixed point sets every observable.}
{Failed. $g_{lag}$ collapses 0/4 unfitted observables over 33 periods (CV $R^2\le0.22$; within-regime $p=0.07$, power 0.89). Observables vary $\times3$--$\times127$. Regime and $\log N$ do better.}
{Do not forecast swarm outcomes from the coupling dial. Plot $N$ against $g_{lag}$, with regime as marker. Scaffold and size carry the signal.}
{0.61}{0.91}

\begin{frame}
\vfill{\Large\bfseries Goals as fields: quench, remanence, memory}\par\vspace{0.8em}
{\small
\chip{ink}{H54}\ The kickoff text sets the day-1 target {\color{muted}(EU 2.60)}\\[0.25em]
\chip{ink}{H96}\ A goal switch is a quench, not hysteresis {\color{muted}(EU 1.11)}\\[0.25em]
\chip{ink}{H82}\ Each agent carries the old goal into the new one {\color{muted}(EU 1.09)}\\[0.25em]
\chip{ink}{H97}\ A kickoff erases position, more along the goal {\color{muted}(EU 1.06)}\\[0.25em]
\chip{ink}{H88}\ A finished goal fades in a week to a small floor {\color{muted}(EU 0.93)}\\[0.25em]
\chip{ink}{H103}\ Content memory fades with work, not with nights {\color{muted}(EU 0.59)}}
\vfill
\end{frame}
% H54 thesis slide. Model tags: card `models` field (role, outcome).
\Models{\mdl{11}{p}{s}\mdl{10}{s}{m}\mdl{15}{s}{u}}
\Math{\mathbf x_i=\mathbf h_{\rm kickoff}+\mathbf b_i+J\,\bar{\mathbf x}+\boldsymbol\xi_i}
\thesis{H54}{Kickoff quench: the day-1 order parameter aligns with the kickoff field}
{H54-kickoff-quench-target/fig.pdf}
{The kickoff is a quench: a field on vector spins.}
{Own kickoff ranks first in 18/33 periods ($p=3\times10^{-6}$). In \#51 each agent lands on its own private goal (0.95).}
{Write the shared target into the kickoff. To move one agent, rewrite its goal text. ``More specific'' does not tighten.}
{0.92}{2.60}

% H96 thesis slide. Model tags: card `models` field (role, outcome).
\Models{\mdl{11}{p}{r}\mdl{15}{t}{s}\mdl{01}{r}{a}}
\Math{R_1=\frac{M(\text{day }1)}{M(\text{pre})},\quad \text{HH: }\ln\tau_{\rm sw}=a+b\,q}
\thesis{H96}{A goal switch is a quench, not hysteresis}
{H96-goal-switch-hysteresis/fig.pdf}
{The old goal state is a magnet. A coercive field would make it linger.}
{Day-1 remanence is 0.27 (bge) / 0.18 (gte), vs 0.85--0.89 across a night. Below the night level in 20/23 transitions. The order test has power 0.50, and its sign is negative.}
{When you reassign a swarm, there is nothing to clear first. Expect stray old-goal references on day 1, not a lag.}
{0.74}{1.11}

% H82 thesis slide. Model tags: card `models` field (role, outcome).
\Models{\mdl{11}{p}{m}\mdl{04}{s}{u}\mdl{12}{t}{a}}
\Math{v_{i,d}=\textstyle\sum_k\alpha_kX_k+\beta\,p_i+\gamma_d\,e_{P-1}+\varepsilon}
\thesis{H82}{Each agent carries the old goal into the new one}
{H82-remanence-endogenous-field/fig.pdf}
{Day-1 content loads on the previous village centroid: an endogenous field $\gamma$.}
{$\Delta\gamma_1=0.16$ / 0.15, $3\times$ the null, 25 boundaries. No decay over 5 days. With the agent's own past as a regressor, the village term drops to $-0.02$ / $-0.06$ (post hoc).}
{Expect about a week of mixed content after a goal change. To stop an old thread, reset that agent's context.}
{0.73}{1.09}

% H97 thesis slide. Model tags: card `models` field (role, outcome).
\Models{\mdl{11}{p}{m}\mdl{02}{s}{m}\mdl{10}{r}{u}}
\Math{v_i^{(1)}=t+\beta_i\,(v_i^{(0)}-t)+\eta_i,\quad \chi_i=1-\beta_i}
\thesis{H97}{A kickoff erases position, more along the goal}
{H97-quench-restoring-force/fig.pdf}
{Each agent relaxes into the kickoff target like a spin in a well.}
{Position memory falls 0.68$\to$0.42 ($\Delta\rho=+0.26\pm0.04$, 16/18 kickoffs). Forgetting: 0.63 along the goal, 0.25 across. Day 1 overshoots. The per-agent rate is not shown stable.}
{Do not read day 1 as the steady state. To retarget one agent, change only its goal text. Measure memory with a disattenuated correlation.}
{0.70}{1.06}

% H88 thesis slide. Model tags: card `models` field (role, outcome).
\Models{\mdl{13}{p}{r}\mdl{04}{s}{u}}
\Math{f(k)=A\,e^{-k/\tau}+c,\quad \tau\approx5.5\ \text{days}}
\thesis{H88}{Collective memory relaxes to a floor in $\approx$1 week; not biexponential}
{H88-collective-memory-decay/fig.pdf}
{Attention to old items decays. Candia predicts a fast and a slow compartment.}
{The biexponential wins 5/27 fits: it fails. Terms: $\tau=5.5$ village days to a 1.9\% floor (20 periods). Newcomers forget at the veterans' pace (ratio 0.93; power 0.87--0.92).}
{To reuse old work, name it again in a kickoff or a file. Before an agent retires, give its live repos to another agent.}
{0.62}{0.93}

% H103 thesis slide. Model tags: card `models` field (role, outcome).
\Models{\mdl{11}{p}{r}\mdl{15}{s}{m}}
\Math{A_w=A_\infty+\Delta A\,f(t),\quad f\in\{\lambda^{N},\,e^{-H/\tau_H}\}}
\thesis{H103}{Remanence decays per unit of work, not per night}
{H103-nights-demagnetize/fig.pdf}
{Kickoff alignment is a remanent field. Test: does it drop one step per night?}
{Active hours fit as well as nights or better in 17/22 periods (16/22 gte). Night step at fixed active lag: $\beta_N=-0.001$ [$-0.023$, 0.021], 41 units. Context resets change nothing.}
{Count content memory in active work hours. An overnight pause does not reset the swarm. To clear a stale theme, change the goal.}
{0.59}{0.59}

\begin{frame}
\vfill{\Large\bfseries Rooms as magnetic domains}\par\vspace{0.8em}
{\small
\chip{ink}{H05}\ Rooms couple chat only through reading {\color{muted}(EU 2.20)}\\[0.25em]
\chip{ink}{H47}\ Rooms bound coherence when they do different work {\color{muted}(EU 1.65)}\\[0.25em]
\chip{ink}{H100}\ Rooms diverge on their own, not along the brief {\color{muted}(EU 1.06)}\\[0.25em]
\chip{ink}{H102}\ Room domains are real; hoppers make no wall {\color{muted}(EU 0.77)}\\[0.25em]
\chip{ink}{H108}\ The room split direction holds; pinning is unclear {\color{muted}(EU 0.21)}}
\vfill
\end{frame}
% H05 thesis slide. Model tags: card `models` field (role, outcome).
\Models{\mdl{02}{p}{m}\mdl{01}{s}{m}\mdl{09}{r}{u}\mdl{15}{t}{a}}
\Math{J=\begin{pmatrix}J_{\rm in}&J_{\rm out}\\ J_{\rm out}&J_{\rm in}\end{pmatrix},\quad \text{cut: } J_{\rm out}\to0}
\thesis{H05}{Rooms couple chat only through reading}
{H05-rooms-cut/fig.pdf}
{Rooms as two coupled spin blocks.}
{Within-room $>$ cross-room talk coupling in 6/6 windows. With read counts in the model, co-location adds nothing. Work ignores rooms.}
{A room is a read filter. To decouple two agents, stop them reading each other. Rooms do not divide work.}
{0.88}{2.20}

% H47 thesis slide. Model tags: card `models` field (role, outcome).
\Models{\mdl{11}{p}{m}}
\Math{C_B=\rho_{\rm cross}/\rho_{\rm within}\le0.3}
\thesis{H47}{Rooms bound coherence when they do different work}
{H47-room-coherence-length/fig.pdf}
{A room boundary cuts the correlation length of content fluctuations.}
{$C_B\le0.3$ in 6/9 periods (median 0.23). Same-task rooms share about half (0.54). The NE42 merge joins rooms; the split restores the bound (DiD 1.00, $p$ 0.001). No room leads.}
{To keep a group's content independent, give it its own room and a different task. Monitor conversation pairs, not room means.}
{0.66}{1.65}

% H100 thesis slide. Model tags: card `models` field (role, outcome).
\Models{\mdl{11}{p}{m}\mdl{10}{s}{m}\mdl{08}{r}{u}}
\Math{\mathbf x_{i,d}=\mathbf g_d+\mathbf a_i+\chi\,\mathbf h_{r}+\mathbf m_{r}+\boldsymbol\epsilon_{i,d}}
\thesis{H100}{Rooms diverge on their own, not along the brief}
{H100-room-symmetry-breaking/fig.pdf}
{Room content = day field + agent constant + operator field $\mathbf h_r$ + self-made part $\mathbf m_r$.}
{After composition and field removal, rooms split beyond random groupings in 5/7 periods (4/7 gte). Composition share $\le$0.18. The split ignores the kickoff text and resets each goal.}
{Expect separate rooms to diverge even with identical instructions. Regroup freely at a goal change. Check a moved agent's work, not its room.}
{0.71}{1.06}

% H102 thesis slide. Model tags: card `models` field (role, outcome).
\Models{\mdl{11}{p}{r}\mdl{01}{s}{m}}
\Math{s_i=\frac{(\mathbf x_i-\mathbf c_{\rm home})\cdot\mathbf u}{(\mathbf c_{\rm other}-\mathbf c_{\rm home})\cdot\mathbf u}\approx\kappa\,p_i}
\thesis{H102}{Room domains exist ($D\approx6$); bridging agents pin no domain wall}
{H102-room-domain-walls/fig.pdf}
{Two rooms are two spin domains. Coupling puts hoppers inside the wall.}
{Rooms with different work form two domains ($D$ 5.3--6.1). The frequent hopper sits at $s=0.06$, inside its home domain. Read dose: $\kappa_R=-0.014$ [$-0.042$, 0.024]. The wall model failed.}
{To keep content apart, give rooms different work. Visitors carry nothing home. To move content, post in the other room or move the agent.}
{0.51}{0.77}

% H108 thesis slide. Model tags: card `models` field (role, outcome).
\Models{\mdl{11}{p}{m}}
\Math{D_\theta=-\ln P(1),\quad R_D=D_\theta^{\rm identical}/D_\theta^{\rm fielded}}
\thesis{H108}{The room split direction holds; pinning is unclear}
{H108-goldstone-room-wandering/fig.pdf}
{Without a field the room direction diffuses (Goldstone). A field pins it.}
{Day-to-day cosine 0.69--0.95 in 6/8 periods. Fielded weeks hold more (0.93 vs 0.70; $R_D=4.7$ [2.1, 5.9]). Without \#38, $R_D=1.4$. Field and split size are not separated.}
{Once rooms split, expect the split direction to hold for the period. Do not treat room instructions as a proven pin.}
{0.42}{0.21}

\begin{frame}
\vfill{\Large\bfseries The context window carries the value}\par\vspace{0.8em}
{\small
\chip{ink}{H44}\ A context wipe costs output, but the agent does not heat up {\color{muted}(EU 3.23)}\\[0.25em]
\chip{ink}{H15}\ Memory loss costs nothing visible; a context wipe costs a tenth of output {\color{muted}(EU 2.73)}\\[0.25em]
\chip{ink}{H69}\ Restatement loops copy the context window {\color{muted}(EU 2.34)}\\[0.25em]
\chip{ink}{H16}\ Idle traps deepen with time {\color{muted}(EU 2.10)}\\[0.25em]
\chip{ink}{H87}\ Only the context window buys commits per bit {\color{muted}(EU 1.57)}\\[0.25em]
\chip{ink}{H71}\ Memory size returns to a set point without overshoot {\color{muted}(EU 1.11)}}
\vfill
\end{frame}
% H44 thesis slide. Model tags: card `models` field (role, outcome).
\Models{\mdl{02}{p}{m}\mdl{04}{s}{s}\mdl{13}{s}{u}}
\Math{P(a)\propto e^{\,h^{0}_a+h^{\rm room}_a+h^{\rm task}_a},\quad \text{wipe: } h^{\rm task}\!\to 0}
\thesis{H44}{A context wipe costs output, but the agent does not heat up}
{H44-context-wipe/fig.pdf}
{The wipe is a field quench, not a temperature pulse.}
{Read share 0.22$\to$0.47 at the first call, relaxes in 5--7 calls. Commits $-38\%$. Entropy \emph{falls}: no temperature pulse. Loops recur 16\% vs 72\%.}
{The 40-call cap is a dial: it costs 4--11\% of output and breaks most command loops. Point a wiped agent at its files.}
{0.88}{3.23}

% H15 thesis slide. Model tags: card `models` field (role, outcome).
\Models{\mdl{04}{p}{m}}
\Math{\kappa_c=\frac{\Delta\mathcal V_c}{I_c}}
\thesis{H15}{Semantic information: memory carries $\approx0$; context carries viability}
{H15-semantic-information-scrambles/figures/summary_obs.pdf}
{Kolchinsky--Wolpert: scramble channel $c$, measure the loss of viability $\mathcal V$.}
{Memory loss: no day-scale cost. Forced wipe: commits $-39\%$ for 10 calls. A longer memory note does not help.}
{Raise the cap or carry task state across the reset. Do not spend effort to protect memory files.}
{0.78}{2.73}

% H69 thesis slide. Model tags: card `models` field (role, outcome).
\Models{\mdl{02}{p}{m}\mdl{08}{s}{s}}
\Math{\mathrm{OR}=\frac{\Pr(\text{copy}\mid\text{in context})}{\Pr(\text{copy}\mid\text{erased})}}
\thesis{H69}{Restatement loops copy the context window}
{H69-restatement-loops/fig.pdf}
{A loop is a fixed point $x_{t+1}\approx x_t$ held by the context.}
{In-context statements are copied $3.4\times$ more than wiped ones (4/4 periods). A wipe: exit $\times2.4$, onset $\times0.54$.}
{Trim the agent's own chat from its context. Nudges do not break loops (0.81).}
{0.85}{2.34}

% H16 thesis slide. Model tags: card `models` field (role, outcome).
\Models{\mdl{02}{p}{r}\mdl{01}{s}{s}\mdl{11}{s}{u}}
\Math{h_{\rm Kramers}(t)=\mathrm{const},\quad h_{\rm Bouchaud}(t)\sim t^{-\mu}}
\thesis{H16}{Idle traps deepen with time}
{H16-trap-aging/fig.pdf}
{A broad spread of trap depths gives aging: escape gets harder with time.}
{Escape slope $-0.77$ in \#51; negative in 6/8 periods. One directed message: odds $\times1.5$--$2.9$. More messages add little.}
{Kick a stuck agent early (hazard $\sim t^{-0.8}$). Send one message. For failure loops, use a reset.}
{0.70}{2.10}

% H87 thesis slide. Model tags: card `models` field (role, outcome).
\Models{\mdl{04}{p}{m}\mdl{12}{t}{a}}
\Math{\kappa_c=\frac{\Delta V_c}{I_c},\quad \Delta V_c=\bar V\,(1-e^{-b_2})}
\thesis{H87}{Value per bit: only the context window has $\kappa_c>0$}
{H87-kappa-table/fig.pdf}
{Each channel has a natural scramble; value per bit $\kappa_c$ ranks the channels.}
{Context: $\kappa_C=5.2$ [3.8, 7.9]; a wipe costs 41\% of output. Own artifact: most bits (0.14), but $\kappa_A=-0.6$ [$-1.6$, 0.3]. Memory, chat, search, human: value consistent with 0.}
{Protect the context window; a longer call cap pays. Notes and search do not buy it back. Test a status message after a wipe.}
{0.63}{1.57}

% H71 thesis slide. Model tags: card `models` field (role, outcome).
\Models{\mdl{04}{p}{m}\mdl{05}{s}{m}}
\Math{x_n-\mu_i=\phi^+\,(x_{n-1}-\mu_i)+\varepsilon_n}
\thesis{H71}{Memory size is an overdamped first-order homeostat}
{H71-memory-homeostat/fig.pdf}
{Each consolidation pulls log memory size back toward the agent's own set point $\mu_i$.}
{$0<\phi^+<1$ in 37/37 units; pooled 0.67 [0.63, 0.71]. No overshoot after forced wipes. Not first order: AR(2) $+0.19$, the target drifts slowly. H09's overshoot was a sampling artifact.}
{Set memory alarms per agent, not as one global cap. Use compression rows only. A prompt change moves the target, not the gain.}
{0.74}{1.11}

\begin{frame}
\vfill{\Large\bfseries Identity: style and family fields}\par\vspace{0.8em}
{\small
\chip{ink}{H46}\ Style is a fingerprint, not a conserved charge {\color{muted}(EU 2.44)}\\[0.25em]
\chip{ink}{H13}\ Model families share a writing style, not a coupling {\color{muted}(EU 2.17)}}
\vfill
\end{frame}
% H46 thesis slide. Model tags: card `models` field (role, outcome).
\Models{\mdl{11}{p}{m}\mdl{04}{s}{r}}
\Math{T=\big\langle\mathrm{rank}\,|\Delta\mathbf s_{\rm bdry}|\big\rangle,\quad T_{\rm conserved}=\tfrac12}
\thesis{H46}{Style is a context-dressed charge: not conserved, still identifying}
{H46-style-conserved-charge/figures/summary_obs.pdf}
{A boundary (goal switch, wipe) is a collision. Does style survive it?}
{Style moves at goal switches and wipes. But it identifies agents at 0.76 vs 0.51 for content. It drifts as context fills.}
{Use the 20 style rates to attribute messages or flag a model swap. Compare at the same context position.}
{0.75}{2.44}

% H13 thesis slide. Model tags: card `models` field (role, outcome).
\Models{\mdl{11}{p}{r}\mdl{02}{s}{r}\mdl{01}{s}{r}\mdl{10}{s}{u}}
\Math{\mathbf x_{i,d}=\mathbf h_d+\mathbf f_{{\rm lab}(i)}+\mathbf b_i+\boldsymbol\xi_{i,d}}
\thesis{H13}{Family fields are a static style offset; no family-block coupling}
{H13-family-fields/figures/summary_obs.pdf}
{Is the family field $\mathbf f$ a position, or a style offset?}
{Family field in 9/15 units; after style removal, 0/9. No family coupling. Room term $17\times$ the lab term.}
{Remove style before any embedding consensus metric. Do not infer vendor ``ideology'' from content.}
{0.87}{2.17}

\begin{frame}
\vfill{\Large\bfseries Collective order: herding, allocation, culture}\par\vspace{0.8em}
{\small
\chip{ink}{H11}\ Agents herd onto shared projects, whatever the goal {\color{muted}(EU 2.06)}\\[0.25em]
\chip{ink}{H81}\ Regime-I culture drifts as a whole over about a month {\color{muted}(EU 1.07)}\\[0.25em]
\chip{ink}{H35}\ The nudger buys glances, not committed work {\color{muted}(EU 1.07)}\\[0.25em]
\chip{ink}{H89}\ Village content changes by transmission, not selection {\color{muted}(EU 0.98)}\\[0.25em]
\chip{ink}{H94}\ Work goes to repos at max entropy, priced by ownership {\color{muted}(EU 0.92)}}
\vfill
\end{frame}
% H11 thesis slide. Model tags: card `models` field (role, outcome).
\Models{\mdl{10}{p}{r}\mdl{13}{s}{u}\mdl{15}{s}{u}}
\Math{P(\boldsymbol\sigma)\propto\exp\big[\textstyle\sum_i h_{\sigma_i}+\frac{\beta J}{N}\sum_{i<j}\delta_{\sigma_i\sigma_j}\big]}
\thesis{H11}{Agents herd onto shared projects, whatever the goal}
{H11-herding-shared-projects/fig.pdf}
{Projects are Potts states. $\beta J>0$ piles agents onto one; $\beta J<0$ spreads them.}
{The sign ignores the goal: $\beta J>0$ in 11/14 weeks. Coupling beyond agent fields in 10/11 shared weeks. Work commits herd too (4/6). Spread only with own artifacts.}
{Expect herding by default. To get parallel work, give each agent its own artifact or role. A coupling $z\ge2$ on artifact links flags duplicated work.}
{0.75}{2.06}

% H81 thesis slide. Model tags: card `models` field (role, outcome).
\Models{\mdl{11}{p}{m}\mdl{14}{t}{a}\mdl{04}{s}{u}}
\Math{\langle u_b\cdot u_{b'}\rangle\propto e^{-|t_b-t_{b'}|/\tau_u}}
\thesis{H81}{Regime-I culture drifts as a whole over about a month}
{H81-culture-slow-mode/fig.pdf}
{Village content has a collective slow mode $u_b$ beyond its members and goal fields.}
{Regime I: contrast 0.24 / 0.21 vs null 0.08 / 0.09; $\tau_u=23$--28 d; it survives turnover. The artifact record does not carry it (power 1.0). Calendar-time drift is not excluded.}
{Compare a village with its own last 2--3 weeks. Curating artifacts or operator wording will not steer this mode.}
{0.61}{1.07}

% H35 thesis slide. Model tags: card `models` field (role, outcome).
\Models{\mdl{04}{p}{m}\mdl{02}{s}{m}}
\Math{\eta_{SU}=\frac{\Delta V}{V^*(I)},\quad \Delta V\le s\sqrt{2rI}}
\thesis{H35}{The nudger is an inefficient Maxwell demon: bits buy activity, not work}
{H35-nudger-demon/fig.pdf}
{The nudger is a Maxwell demon: it spends bits of agent state to buy activity.}
{\#51: 1.4 bits per nudge buy $+1.45$ active min, at 38\% of the frontier. Commits $+0.22$/h [$-0.29$, 0.82], not significant. Switching the nudger on or off does not move swarm work.}
{Nudge each pause chain once, early ($k=2$--3): $\times2.3$ active minutes. Stop re-nudging deep traps. Do not expect more output.}
{0.61}{1.07}

% H89 thesis slide. Model tags: card `models` field (role, outcome).
\Models{\mdl{13}{p}{m}\mdl{05}{s}{r}\mdl{06}{n}{a}}
\Math{\Delta\bar z=\mathrm{Cov}_S(w,z)+\mathrm{E}_S(w\,\Delta z)+\mathrm{Mig}}
\thesis{H89}{Price decomposition: culture changes by transmission; selection $\approx0$}
{H89-price-equation/fig.pdf}
{Price equation: change of the mean trait = selection + transmission + migration.}
{Transmission carries 0.98 of the change (median, 33/33 periods). Selection 0.01: the design cannot detect it. Style migration exceeds content in 21/26. No attractor ($C>0$ in 0/29).}
{Read a content shift as a day field on all agents. A style shift on a join or exit day is mostly composition.}
{0.78}{0.98}

% H94 thesis slide. Model tags: card `models` field (role, outcome).
\Models{\mdl{10}{p}{s}\mdl{05}{s}{m}\mdl{06}{r}{m}}
\Math{P_{ij}\propto a_i\,b_j\,e^{\lambda_{\rm own}{\rm own}_{ij}+\lambda_{\rm room}{\rm room}_{ij}}}
\thesis{H94}{Work allocation is a max-entropy state; ownership is the Lagrange multiplier}
{H94-maxent-work-allocation/fig.pdf}
{Work is a max-ent agent--repo table; each constraint acts as a price $\lambda$.}
{Residual $\le0.13$ of $I(\text{agent};\text{repo})$ in 23/24 units. Ownership price 2.2 nats in shared weeks, 7.8 in own-role units ($p=0.0007$). No specialization beyond ownership.}
{Track $\lambda_{\rm own}$ and $\kappa$ each week. To spread work, assign owners. To pool it, name one shared repo in the kickoff.}
{0.61}{0.92}



\end{document}
